\section{Propagation to time zero and localization}\label{sec:contact-localization}

We now apply the discrete inequalities to the extremal pair $v,W$ and the
constant $\kappa$ supplied by Proposition~\ref{prop:envelope-contact}.  Write
$z_*=(0,1)$ and $H_*=H_v(z_*)\in\R$; the starting limit is
\eqref{eq:dp-J-contact}.

\begin{proposition}\label{prop:contact-propagation}
Put $\delta_j=2^{-j}$ for $j\ge1$.  There exist a constant $K<\infty$, a point
$x_{\rm end}\in\R^m$, and an infinite set $\mathcal J\subset\mathbb N$ with the
following property.  For every $j\in\mathcal J$ there are a point
$x_j\in\R^m$ and a fit $(p_j,a_j)$ for $v$ at $(x_j,\delta_j)$ of radius
$\delta_j$, in the sense of \eqref{eq:limit-fit}, such that
\begin{equation}\label{eq:contact-propagation-conclusion}
\begin{gathered}
  H_v(x_{\rm end},0)>-\infty,\qquad
  W(x_j,\delta_j)=v(x_j,\delta_j),\\
  |x_{\rm end}-x_j+p_j\delta_j|\le K\delta_j.
\end{gathered}
\end{equation}
\end{proposition}

The last bound says that the same terminal point stays in a fixed bounded
part of every selected fitted frame, where its finite slope will identify
the limiting fixed-length slopes. Figure~\ref{fig:contact-localization}
shows how two selected fitted frames localize the same endpoint.

\begin{figure}[tbp]
\centering
\begin{tikzpicture}[x=1cm,y=1cm,>=Latex,font=\small]
  \definecolor{fitblue}{RGB}{35,88,136}
  \definecolor{fitorange}{RGB}{168,83,31}
  \fill[fitblue!9] (-1.00,0)--(2.50,0)--(3.25,3.15)--(-0.25,3.15)--cycle;
  \draw[fitblue!50] (-1.00,0)--(-0.25,3.15) (2.50,0)--(3.25,3.15);
  \draw[fitblue!50] (-0.25,3.15)--(3.25,3.15);
  \fill[fitorange!12] (-1.225,0)--(0.525,0)--(-0.045,1.575)--(-1.795,1.575)--cycle;
  \draw[fitorange!65] (-1.225,0)--(-1.795,1.575) (0.525,0)--(-0.045,1.575);
  \draw[fitorange!65] (-1.795,1.575)--(-0.045,1.575);
  \draw[fitblue,densely dashed,line width=0.85pt] (0.75,0)--(1.50,3.15);
  \draw[fitorange,dash pattern=on 4pt off 2pt on 1pt off 2pt,line width=0.85pt]
      (-0.35,0)--(-0.92,1.575);
  \draw[->] (-2.65,0)--(3.65,0) node[right] {$x$};
  \draw[->] (-2.45,-0.02)--(-2.45,3.80) node[above] {$t$};
  \draw[black!35,dotted] (-2.45,3.15)--(1.50,3.15);
  \draw[black!35,dotted] (-2.45,1.575)--(-0.92,1.575);
  \node[left] at (-2.48,3.15) {$\delta_j$};
  \node[left] at (-2.48,1.575) {$\delta_k$};
  \node[left] at (-2.48,0) {$0$};
  \fill[fitblue] (1.50,3.15) circle (2.3pt);
  \node[above=5pt,fitblue] at (1.50,3.15) {$(x_j,\delta_j)$};
  \fill[fitorange] (-0.92,1.575) circle (2.3pt);
  \node[above left=3pt,fitorange] at (-0.92,1.575) {$(x_k,\delta_k)$};
  \fill (0,0) circle (2.8pt);
  \node[below=5pt] at (0,0) {$(x_{\rm end},0)$};
  \draw[fitblue] (0.75,-0.045)--(0.75,0.045);
  \draw[fitorange] (-0.35,-0.045)--(-0.35,0.045);
  \node[anchor=west,align=left,fitblue] (ja) at (4.00,2.80)
      {$x=x_j+p_j(t-\delta_j)$};
  \draw[fitblue!75] (ja.west)--(3.42,2.80)--(1.29,2.28);
  \node[anchor=west,align=left,fitorange] (ka) at (4.00,1.68)
      {$x=x_k+p_k(t-\delta_k)$};
  \draw[fitorange!75] (ka.west)--(3.42,1.68)--(-0.72,1.13);
  \node[anchor=west,align=left,text width=4.25cm] at (4.00,0.46)
      {$j,k\in\mathcal J$, $j<k$.\\[3pt]
       One endpoint lies in both strips.};
  \draw[fitblue,<->] (0.75,-0.72)--(2.50,-0.72)
      node[midway,below=2pt] {$K\delta_j$};
  \draw[fitblue!50] (0.75,-0.64)--(0.75,-0.82)
                        (2.50,-0.64)--(2.50,-0.82);
\end{tikzpicture}
\caption{Schematic localization at two selected scales, shown in one spatial
 direction with different horizontal and vertical units. The dashed axes
 belong to the two fitted frames; each shaded endpoint-localization strip has spatial half-width
 $K\delta_j$ or $K\delta_k$. Their time-zero sections both contain the
 same point $(x_{\rm end},0)$, expressing
 $|x_{\rm end}-x_j+p_j\delta_j|\le K\delta_j$.
 The conclusion holds for infinitely many selected indices
 $j\in\mathcal J$. The axes are separate affine fits, not a trajectory
 joining the marks.}
\label{fig:contact-localization}
\end{figure}


\begin{proof}
We first construct finite stopped walks with common expectation and error
bounds, and then localize their terminal points in the frames of their
marks.  Finally, one joint limiting law will retain the terminal point and
all marks, allowing us to select infinitely many localized contacts for the
same endpoint.

\medskip\noindent\emph{The stopped walks and their expectation bounds.}
Fix for the moment an integer $T\ge2$ and a radius $R\ge2$.  For each
$N\ge1$ choose $\epsilon_N$ to be an inverse power of two with
\begin{equation}\label{eq:contact-grid}
  0<\epsilon_N\le\frac{\delta_N}{4}.
\end{equation}
In particular, $\epsilon_N$ divides $1,T$, and every $\delta_j$ with
$j\le N$.  We will make each $\epsilon_N$ smaller if necessary, retaining
these conditions.  Write $\epsilon=\epsilon_N$ while constructing the
$N$-th walk.

Start at $z_0=z_*$.  At a state $z_n=(x_n,t_n)$ with
$\delta_N<t_n<T$ and $|x_n|<R$, choose the two steps of
\eqref{eq:dp-step-choice} and take $z_{n+1}=z_n^+$ or $z_n^-$ with equal
probabilities.  Stop on first reaching $t_n=\delta_N$, $t_n=T$, or
$|x_n|\ge R$; call this stopping index $\tau_N$.  The choices of optimizers
can be made separately at every node of the countable tree of finite
histories.  Taking independent fair signs $\xi_n\in\{-1,1\}$ to select the
steps then defines the walk and the filtration
$\mathcal F_n=\sigma(\xi_0,\ldots,\xi_{n-1})$; all velocities at step $n$
are $\mathcal F_n$-measurable.  In expressions such as $Q_-(z_n)$ below,
the velocity denotes the choice made at that history.  Before stopping,
$t_n>\delta_N\ge4\epsilon$, so the inequality \eqref{eq:dp-drift} applies.

We choose $\epsilon_N$ sufficiently small that every position up to and
including the stop lies in $\overline B_{R+1}(0)$.  Here is the uniform
localization of optimizers that permits this choice.  The sandwich for $W$
extends it continuously to $t=0$ by $W(x,0)=v(x,0)$.  The growth bounds for
$v$ and the sandwich give uniform subquadratic growth for $v$ and $W$ on
compact time intervals in $\R$ and $[0,\infty)$, respectively.  For base
points in a fixed compact box and lengths
$\ell\le1$, a quadratic cost first confines all optimizing positions for
$\Penv_\ell W$ and $\Menv_\ell v$ to a fixed spatial compact set: outside a
sufficiently large ball the cost dominates the uniform subquadratic bound,
even with the weaker coefficient corresponding to $\ell=1$.  On that compact
set the functions are bounded.  Comparing an optimizer with the zero
displacement competitor then gives $|y-x|^2/(2\ell)\le C$ there, uniformly
as $\ell\downarrow0$.  Thus the optimizing displacements tend to zero
uniformly.  Apply this to length $2\epsilon$ and base points with
$|x|\le R$, $0\le t\le T$, in a slightly larger time interval.  The
intermediate displacements used for the steps are half the full
displacements, so they can all be made at most one.  This proves the asserted
bound at stopping.  We also choose $\epsilon_N$ small enough that
\begin{equation}\label{eq:contact-start-lower}
  J_{\epsilon_N}(z_*)\ge H_*-1,
\end{equation}
which is possible by \eqref{eq:dp-J-contact}.

The time coordinate before stopping is the simple symmetric walk with step
$\epsilon$ starting at $1$.  It is stopped no later than this time walk hits
$\delta_N$ or $T$.  The latter hitting time is almost surely finite and has
finite expectation, for example by the elementary gambler's ruin formula.
The same is therefore true of $\tau_N$.  Put
\[
  \mathcal B=\overline B_{R+1}(0)\times[0,T],\qquad B=R+T+2.
\]
The upper slope bound for $v$ on boxes and \eqref{eq:dp-J-bounds} give a
constant $U_B=C L(B)^4$, independent of $N$, such that
\begin{equation}\label{eq:contact-J-upper}
  J_{\epsilon_N}(z_n)\le U_B\qquad (0\le n\le\tau_N).
\end{equation}

For any stopping time $\sigma\le\tau_N$, summing the conditional form of
\eqref{eq:dp-drift} up to $\sigma\wedge h$ gives
\begin{equation}\label{eq:contact-truncated-drift}
\begin{split}
  \mathbb E J_\epsilon(z_{\sigma\wedge h})
  &\ge J_\epsilon(z_*)\\
  &\quad+c\,\mathbb E\sum_{n<\sigma\wedge h}
          \bigl(|Q_+(z_n)-Q_0(z_n)|^2+|Q_-(z_n)-Q_0(z_n)|^2\bigr).
\end{split}
\end{equation}
For each fixed $N$, $J_\epsilon$ is continuous and bounded on the compact
box $\overline B_{R+1}(0)\times[\delta_N,T]$ containing the stopped walk.
Thus dominated convergence applies to its
term as $h\to\infty$.  Monotone convergence applies to the nonnegative sum.
We obtain, for all these stopping times,
\begin{equation}\label{eq:contact-stopped-expectation}
  \mathbb E J_\epsilon(z_\sigma)\ge J_\epsilon(z_*).
\end{equation}
Taking $\sigma=\tau_N$ in \eqref{eq:contact-truncated-drift} and using
\eqref{eq:contact-start-lower}--\eqref{eq:contact-J-upper} also gives
\begin{equation}\label{eq:contact-error-budget}
  \mathbb E\sum_{n<\tau_N}
       \bigl(|Q_+(z_n)-Q_0(z_n)|^2+|Q_-(z_n)-Q_0(z_n)|^2\bigr)
  \le B_0,\qquad B_0=C_*L(B)^4.
\end{equation}
Here $C_*$ is fixed for $v$ and independent of $N,R,T$; it absorbs the
finite number $H_*$ and the absolute constant $c$ in \eqref{eq:dp-drift}.
Both \eqref{eq:contact-stopped-expectation} and
\eqref{eq:contact-error-budget} are unconditional expectations. Together
with the upper slope bound, the lower expectation bound will give contact
and a finite terminal slope in the limit. We first use the discrepancy
bound to control horizontal-exit probabilities, both in the original
coordinates and in fitted frames.

\medskip\noindent\emph{Reaching the lower level with positive probability.}
We next show how to choose the outer cutoffs so that the walks reach
$t=\delta_N$ with probability at least $3/4$, uniformly in $N$.
Couple the time signs with
the simple symmetric walk starting at $1$ and stopped only at $0$ or $T$.
Its probability of reaching $T$ is $1/T$, which also bounds that probability
for our walk, since it may stop sooner. To control horizontal exit, we will
bound the expected maximal displacement by $C_T L(R+T+2)^2$, with $C_T$ independent
of $R$ and $N$. Markov's inequality will then make the probability of
horizontal exit tend to zero as $R\to\infty$ for fixed $T$.

The expected number of visits of the unrestricted time walk to any fixed
interior grid level is at most $CT/\epsilon$. To verify this directly,
write $q=T/\epsilon$ and identify the level with $k\in\{1,\ldots,q-1\}$.
After a visit to $k$, the
probability of reaching an endpoint before the next return is
\[
  \frac12\left(\frac1k+\frac1{q-k}\right)\ge\frac2q.
\]
This follows from the linear hitting probability for a symmetric walk on
each of the intervals $[0,k]$ and $[k,q]$. Starting from the first visit,
the expected total number of visits is at most $q/2$, by the geometric
return calculation; the probability of that first visit is at most one.
Stopping the actual walk earlier can only decrease its visit counts.
Summing this bound over levels gives
\begin{equation}\label{eq:contact-time-and-velocity-budget}
  \mathbb E(\epsilon^2\tau_N)\le CT^2,\qquad
  \mathbb E\sum_{n<\tau_N}\epsilon^2|Q_0(z_n)|^2
       \le C_T L(B)^2.
\end{equation}
For the second inequality we used Lemma~\ref{lem:dp-log-velocity} before
each step.  More explicitly, it bounds the summand velocity squared by
$C(L(B)^2+|\log t_n|^2)$, while the visit estimate gives at most
\[
  CT\epsilon\sum_{k=1}^{q-1}
           \bigl(L(B)^2+|\log(k\epsilon)|^2\bigr).
\]
This has the asserted bound.  Indeed, on $(0,1]$ the function
$|\log s|^2$ is decreasing and integrable, so its right-endpoint grid sum
times $\epsilon$ is at most its integral there; on $[1,T]$ it is bounded by
$(\log T)^2$.  The constants $C_T$ in this proof can depend on $T$ and the
fixed data, but not on $R$ or $N$.

For $n<\tau_N$ the spatial increment can be written exactly as
\[
  x_{n+1}-x_n
   =\epsilon\xi_n Q_0(z_n)
     +\epsilon\xi_n\bigl(Q_{\xi_n}(z_n)-Q_0(z_n)\bigr),
\]
where $Q_{1}=Q_+$ and $Q_{-1}=Q_-$.  Extend both kinds of increments by
zero after stopping.  The first summands are martingale differences, since
the active-step indicator and $Q_0(z_n)$ are known before the fair sign.
For a deterministic truncation in the number of steps, Doob's $L^2$ maximal
inequality and
orthogonality of these differences give an upper bound
\[
  2\left(\mathbb E\sum_{n<\tau_N}\epsilon^2|Q_0(z_n)|^2\right)^{1/2}
\]
for the expected supremum of the norm of their partial sums.  Letting the
truncation increase preserves this bound by monotone convergence of the
suprema.  On each path, the sum of the norms of the other increments is at
most
\[
  (\epsilon^2\tau_N)^{1/2}
  \left(\sum_{n<\tau_N}
       \bigl(|Q_+(z_n)-Q_0(z_n)|^2+|Q_-(z_n)-Q_0(z_n)|^2\bigr)\right)^{1/2}.
\]
Cauchy--Schwarz in expectation gives the travel inequality
\begin{equation}\label{eq:contact-travel-inequality}
\begin{split}
  \mathbb E\sup_{0\le n\le\tau_N}|x_n|
  &\le 2\left(\mathbb E\sum_{n<\tau_N}
                    \epsilon^2|Q_0(z_n)|^2\right)^{1/2}\\
  &\quad+\left(\mathbb E(\epsilon^2\tau_N)\right)^{1/2}\\
  &\qquad\times\left(\mathbb E\sum_{n<\tau_N}
       \bigl(|Q_+(z_n)-Q_0(z_n)|^2+|Q_-(z_n)-Q_0(z_n)|^2\bigr)\right)^{1/2}.
\end{split}
\end{equation}
This calculation uses only predictability of the coefficients of the fair
signs in the reference increments, including the active-step indicators,
and the pathwise bound on error lengths.  It therefore applies as well to
the displacement from the starting point of a stopped portion after a stopping
time, with its increments extended by zero outside that portion.  Equations \eqref{eq:contact-error-budget} and
\eqref{eq:contact-time-and-velocity-budget} now show that
\begin{equation}\label{eq:contact-original-travel}
  \mathbb E\sup_{0\le n\le\tau_N}|x_n|
    \le C_T\bigl(L(B)+\sqrt{B_0}\bigr)
    \le C'_T L(B)^2.
\end{equation}

The probability of a horizontal stop is at most the right side of
\eqref{eq:contact-original-travel} divided by $R$. For fixed $T$ this tends
to zero as $R\to\infty$. Choose an integer $T\ge8$, and then choose $R$
large enough that
the horizontal-stop probability is at most $1/8$, uniformly in $N$.
Fix these $T,R$ for the rest of the proof, so $B$ and $B_0$ are now fixed
finite constants.  If
\[
  \eta_N=\mathbf1_{\{t_{\tau_N}=\delta_N\}},
\]
the two preceding bounds and the definition of the stop imply
\begin{equation}\label{eq:contact-lower-hit-probability}
  \mathbb P(\eta_N=1)\ge\frac34\qquad(N\ge1).
\end{equation}
This remains true if the lower time level and the horizontal cutoff are
reached simultaneously, since that case has $\eta_N=1$.

\medskip\noindent\emph{Marks and localization in fitted frames.}
For $1\le j\le N$, let
\[
  \sigma_{j,N}=\tau_N\wedge\inf\{n\ge0:t_n=\delta_j\},\qquad
  z_j=(x_j,t_j)=z_{\sigma_{j,N}}.
\]
Thus a mark is the first point on its time level, or the terminal point if
the walk stops before that level.  The divisibility of the grid ensures
$t_j\ge\delta_j$, with equality on $\{\eta_N=1\}$.  At the time of this
mark choose a fit $(p_j,a_j)$ at $z_j$ of radius $\delta_j$.  These choices
can again be made on the countable tree of histories, so the parameters are
$\mathcal F_{\sigma_{j,N}}$-measurable.  The bound on fit parameters on boxes
gives, for a fixed constant independent of $N$,
\[
  |p_j|+|a_j|\le\Gamma_j:=C\bigl(L(B)+j\bigr),\qquad j\le N,
\]
increasing $C$ to make every $\Gamma_j\ge1$.

The cutoffs $T,R$ and the error bound $B_0$ are now fixed.  We will choose
two further cutoffs, uniformly in $j,N$, to control the terminal point in
each fitted frame.  Fix $1\le j\le N$, put $r=\delta_j$, and let
$\mathcal R_{j,N}=\{t_j=r\}$ be the event that the level $r$ is reached,
including the case where it is reached at the original stop.  This event is
$\mathcal F_{\sigma_{j,N}}$-measurable.  On it $z_j=(x_j,r)$, so we may apply
Lemma~\ref{lem:fitted-frame-covariance} pathwise with the chosen fit
$(p_j,a_j)$.  In its coordinates,
\[
  s=\frac{t}{r},\qquad
  y=\frac{x-x_j-p_j(t-r)}{r},\qquad
  \epsilon'=\frac{\epsilon_N}{r}.
\]
The transformed solution $\bar v$ is centered at $(0,1)$ and has the same
Jensen bound.  Its $\bar W$ therefore has the standing properties of
Section~\ref{sec:discrete-propagation}.  The mapped steps are optimizing
steps, and \eqref{eq:contact-velocity-covariance} gives identical squared
velocity discrepancies.  Contact at this finite mark is not needed.  The
fit and frame are known at the marking time and will be held fixed afterward.

On $\mathcal R_{j,N}$ consider the portion of the original walk from its
$j$-th mark to its stop, with the additional rule that this portion stops
on first reaching $s=M$ or $|y|\ge K$, where $M\ge2$ is an integer and
$K\ge2$.  It includes the endpoint of this additional stopping rule.  If
the original walk stopped at the mark, the portion has no steps.  A hit of
an additional cutoff simultaneous with the original stop is counted as a
hit of that additional cutoff.  On $\mathcal R_{j,N}^{\rm c}$ define the
portion to be empty as well, and assign zero to all its sums and spatial
suprema.  Let $n'$ be its number of steps, and write $\sum'$ for a sum over
its pre-step states.

The total squared velocity error on this portion, in the new frame, has the
unconditional bound
\begin{equation}\label{eq:contact-frame-error-budget}
  \mathbb E\sum'
       \bigl(|Q'_+-Q'_0|^2+|Q'_--Q'_0|^2\bigr)\le B_0.
\end{equation}
Indeed, on each path the sum is a subsum of the original sum in
\eqref{eq:contact-error-budget}, with identical summands by
\eqref{eq:contact-velocity-covariance}; on an empty portion it is zero.
This argument gives an unconditional estimate.  It makes no assertion that
the same bound holds after conditioning on a particular history at the
mark.

The time-walk estimates, in contrast, do hold conditionally on that history.
On $\mathcal R_{j,N}$ the new time starts at $s=1$ and has step
$\epsilon'=\epsilon_N/r\le1/4$, which divides $1$ and $M$.  Given
$\mathcal F_{\sigma_{j,N}}$, the subsequent signs remain fair.  The portion
is stopped no later than the unrestricted new time walk first hits $0$ or
$M$.  Its conditional expected visits to each interior level are therefore
at most $CM/\epsilon'$ by the same visit argument used above.  Before any
step of the portion, the original walk is still live, so
$s=t/r>\delta_N/r\ge4\epsilon'$.  Also $|y|<K$ and $s<M$ before the
additional stop.  Lemma~\ref{lem:dp-log-velocity} applies to the centered
$\bar v$ and its $\bar W$ uniformly, and bounds $|Q'_0|$ by
$C(L(K+M+2)+|\log s|)$.  The integrable logarithmic grid sum gives
\begin{equation}\label{eq:contact-frame-time-budget}
\begin{split}
  \mathbb E\bigl((\epsilon')^2n'\mid\mathcal F_{\sigma_{j,N}}\bigr)
       &\le CM^2\mathbf1_{\mathcal R_{j,N}},\\
  \mathbb E\left(\sum'
          (\epsilon')^2|Q'_0|^2\,\middle|\,\mathcal F_{\sigma_{j,N}}\right)
       &\le C_M L(K+M+2)^2\mathbf1_{\mathcal R_{j,N}}.
\end{split}
\end{equation}
The frame is known at the conditioning time, and the estimates of the
centered class have constants independent of that frame.  This explains the
uniformity of these conditional bounds.  Integrating them gives the same
bounds unconditionally without the indicator.

Extend the new increments by zero outside the portion.  Their reference
part is $\epsilon'\xi_nQ'_0$ and their error part is
$\epsilon'\xi_n(Q'_{\xi_n}-Q'_0)$.  The reference increments are martingale
differences after the random starting mark: the fit and reached event are
known at that mark, and on an active step the indicator and
$Q'_0=Q_0-p_j$ are $\mathcal F_n$-measurable while $\xi_n$ is still fair
given $\mathcal F_n$.  Thus the calculation of
\eqref{eq:contact-travel-inequality} applies in the primed variables.  Using
the integrated bounds from
\eqref{eq:contact-frame-time-budget} and the unconditional bound
\eqref{eq:contact-frame-error-budget}, it gives
\begin{equation}\label{eq:contact-frame-travel}
  \mathbb E\left(\sup_{\text{points of the portion}}|y|\right)
   \le C'_M\bigl(L(K+M+2)+\sqrt{B_0}\bigr).
\end{equation}
The supremum is understood to be zero for an empty portion.  In particular,
the probability that this portion reaches $|y|\ge K$ is at most the right
side of \eqref{eq:contact-frame-travel} divided by $K$.  The probability it
reaches $s=M$ is at most $1/M$: conditionally on reaching the mark, the
unrestricted symmetric time walk starting from $1$ hits $M$ before $0$ with
probability $1/M$, and an earlier stop only removes this event.

Choose a fixed integer $M\ge8$, and then choose $K$ so large that
\[
  \frac{C'_M\bigl(L(K+M+2)+\sqrt{B_0}\bigr)}{K}\le\frac18.
\]
This is possible because $B_0$ is already fixed and $L(K+M+2)/K\to0$.
On $\{\eta_N=1\}$ the $j$-th level is reached.  If neither additional
cutoff intervenes, the terminal point is in the portion and has $|y|\le K$.
Combining this observation with \eqref{eq:contact-lower-hit-probability}
and the two additional-cutoff probabilities proves, uniformly for $j\le N$,
\begin{equation}\label{eq:contact-finite-localization}
  \mathbb P\left(\eta_N=1,\quad
      |x_{\tau_N}-x_j-p_j(t_{\tau_N}-\delta_j)|\le K\delta_j\right)
       \ge\frac34-\frac18-\frac18=\frac12.
\end{equation}

\medskip\noindent\emph{One limiting law and the contact properties.}
We use one fixed compact space for all these laws.  Define
\[
  \mathcal K_j=\mathcal B\times\overline B_{\Gamma_j}(0)
                    \times[-\Gamma_j,\Gamma_j],
  \qquad
  \mathcal X=\mathcal B\times\{0,1\}\times\prod_{j=1}^{\infty}\mathcal K_j.
\]
The coordinates are the terminal point, the flag, and the marked point and
fit at each index.  This countable product of compact metric spaces is
compact and metrizable.  For indices $j>N$, fill the coordinates by any
fixed point of $\mathcal K_j$; these dummy entries need not be fits.  Let
$\mu_N$ be the law on $\mathcal X$ of the terminal point $z_{\tau_N}$, the
flag $\eta_N$, and the marked points and fits.  Compactness of the space
of probability measures on a compact metric space gives a subsequence
converging weakly to a probability law $\mu$; all subsequent limits in $N$
are along this subsequence.  Denote its coordinates by
$z_{\rm end}=(x_{\rm end},t_{\rm end})$, $\eta$, and
$(z_j,p_j,a_j)$.

Here are first the properties that pass directly to this law.  The flag is
a coordinate in a discrete space, so its probability converges and
$\mu(\eta=1)\ge3/4$.  Also the continuous nonnegative function
$\eta t_{\rm end}$ on $\mathcal X$ has integral at most $\delta_N$ under
$\mu_N$.  Its integral under $\mu$ is zero.  It follows that
$t_{\rm end}=0$ almost surely on $\{\eta=1\}$.  For each fixed $j$, the sets
$\{t_j\ge\delta_j\}$ and $\{\eta=0\}\cup\{t_j=\delta_j\}$ are closed and
have $\mu_N$-measure one for all $N\ge j$.  The closed-set direction of the
Portmanteau Theorem, $\limsup_N\mu_N(F)\le\mu(F)$, makes their
$\mu$-measures one.  Fits at the specified radius form a closed subset of
$\mathcal K_j$.  To see this without varying the sheared domain, write its
increments in \eqref{eq:limit-fit} as
\[
  \ell=\delta_j\theta,\qquad
  h=p_j\delta_j\theta+\delta_j\zeta,
  \qquad |\theta|\le1,\quad |\zeta|\le1.
\]
For each $(\theta,\zeta)$ the fitting inequality is closed in
$(z_j,p_j,a_j)$ by continuity of $v$, and their intersection over this
fixed compact set of increments is closed.  Portmanteau therefore also
passes the fitting property to $\mu$.  Taking the intersection of the
resulting full-measure sets for all $j$, we have
\begin{equation}\label{eq:contact-mark-limit}
\begin{gathered}
  t_{\rm end}=0\text{ on }\{\eta=1\},\qquad
  t_j\ge\delta_j\text{ and }t_j=\delta_j\text{ on }\{\eta=1\},\\
  (p_j,a_j)\text{ is a fit at }z_j\text{ of radius }\delta_j
       \quad\text{for every }j,
\end{gathered}
\end{equation}
almost surely.

Two further properties come from the expectation bounds:
\begin{equation}\label{eq:contact-limit-properties}
  H_v(z_{\rm end})>-\infty,\qquad W(z_j)=v(z_j)\quad(j\ge1)
  \qquad\text{almost surely under }\mu.
\end{equation}
For the first, fix $\ell>0$ and set
\[
  \Phi_\ell(z)=\frac12\bigl(\Senv^+_\ell v(z)+\Senv^-_\ell v(z)\bigr).
\]
This is a finite continuous function on $\mathcal B$.  For all sufficiently
large $N$, $\epsilon_N\le\ell$, and monotonicity of both slopes with respect
to their length and \eqref{eq:dp-J-bounds} give
$J_{\epsilon_N}\le\Phi_\ell$.  Apply
\eqref{eq:contact-stopped-expectation} at $\tau_N$ and pass to the weak limit
using this fixed continuous function and \eqref{eq:dp-J-contact}.  The
result is
\begin{equation}\label{eq:contact-terminal-slopes}
  \int\Phi_\ell(z_{\rm end})\,d\mu\ge H_*\qquad(\ell>0).
\end{equation}
Take a decreasing sequence $\ell_h\downarrow0$ with $\ell_h\le1$.  The
functions $\Phi_{\ell_h}$ decrease to $H_v$ pointwise, because the two slope
limits agree for $v$.  Their uniform upper bound on $\mathcal B$ is $U_B$
after increasing its fixed constant.  Hence
$U_B-\Phi_{\ell_h}(z_{\rm end})$ are nonnegative and increase to
$U_B-H_v(z_{\rm end})$.  Monotone convergence and
\eqref{eq:contact-terminal-slopes} give
\[
  \int\bigl(U_B-H_v(z_{\rm end})\bigr)\,d\mu\le U_B-H_*<\infty.
\]
The integrand would be infinite where $H_v=-\infty$, proving the first
claim of \eqref{eq:contact-limit-properties}.

For the contact claim fix $j$ and let $\nu_N^j$ and $\nu^j$ be the marginals
of $\mu_N$ and $\mu$ for the marked point $z_j$.  They converge weakly.
Suppose that $\nu^j$ gives positive mass to $\{v-W>0\}$.  All these limit
marks have $t\ge\delta_j$.  Since $v-W$ is continuous and nonnegative for
$t>0$, there is $\alpha>0$ such that the relatively open set
\[
  U=\{(x,t)\in\mathcal B:t>\delta_j/2,\ v(x,t)-W(x,t)>\alpha\}
\]
has $\nu^j(U)>0$ and its closure is contained in
$\{t\ge\delta_j/2,\ v-W\ge\alpha\}$.  For example, first choose $\alpha$
with positive mass where $v-W>2\alpha$.  The zero-length continuity of the
quadratic envelopes is uniform on this compact closure.  Thus
\eqref{eq:dp-contact-gap} gives
\[
  2\epsilon_NJ_{\epsilon_N}(z)
   \le-\frac{\alpha}{2}
\]
on $U$ for all sufficiently large $N$.  In particular,
$J_{\epsilon_N}\le-\alpha/(4\epsilon_N)$ on $U$.  By the open-set direction
of Portmanteau, $\nu^j(U)\le\liminf_N\nu_N^j(U)$, so these probabilities
are bounded below by a fixed positive number for large $N$.  The uniform
upper bound \eqref{eq:contact-J-upper} on all marks and
\eqref{eq:contact-stopped-expectation} at $\sigma_{j,N}$ would then give
\[
  J_{\epsilon_N}(z_*)\le\int J_{\epsilon_N}\,d\nu_N^j
  \le-\frac{\alpha}{4\epsilon_N}\nu_N^j(U)
      +U_B\bigl(1-\nu_N^j(U)\bigr)\longrightarrow-\infty.
\]
This contradicts \eqref{eq:dp-J-contact}.  Thus $W(z_j)=v(z_j)$ almost
surely.  A countable intersection proves this simultaneously for all $j$.
The derivation of \eqref{eq:contact-limit-properties} used unconditional
expectations; it does not require conditioning on the flag.

For each fixed $j$ define the subset of $\mathcal X$
\[
  E_j=\left\{\eta=1,\quad
      |x_{\rm end}-x_j-p_j(t_{\rm end}-\delta_j)|\le K\delta_j\right\}.
\]
This event retains the actual terminal time in the prelimit laws.
It is closed: the flag condition is closed and the norm is a continuous
function of the terminal point, the $j$-th marked point, and $p_j$.
Equation~\eqref{eq:contact-finite-localization} gives
$\mu_N(E_j)\ge1/2$ whenever $N\ge j$.  The
closed-set Portmanteau inequality has the direction
\[
  \mu(E_j)\ge\limsup_{N\to\infty}\mu_N(E_j)\ge\frac12.
\]
These events need not be independent.  For every $h$ their tail union has
measure at least $1/2$, since it contains each $E_j$ with $j\ge h$.  Continuity of
probability from above therefore gives
\[
  \mu\left(\limsup_{j\to\infty}E_j\right)
   =\lim_{h\to\infty}\mu\left(\bigcup_{j\ge h}E_j\right)\ge\frac12.
\]
Choose one outcome in this limsup event and in the full-measure sets of
\eqref{eq:contact-mark-limit}--\eqref{eq:contact-limit-properties}.  Every
$E_j$ includes $\eta=1$, so this outcome has $t_{\rm end}=0$ and
$t_j=\delta_j$ for every $j$.  It belongs to $E_j$ for an infinite set of
indices, at each of which its localization inequality is
\[
  |x_{\rm end}-x_j+p_j\delta_j|\le K\delta_j.
\]
The same outcome has a finite common slope at the endpoint, contact at all
the marks, and the indicated fits.  Its endpoint and the infinite set of
these marks establish \eqref{eq:contact-propagation-conclusion}.
\end{proof}
